\section*{Capítulo \ref{ch:g10:electron-shells} --- Las capas electrónicas y la tabla periódica}

\begin{solution}{exo:g10:electron-shells:1}
C: $1\mathrm{s}^2\,2\mathrm{s}^2\,2\mathrm{p}^2$. N:
$1\mathrm{s}^2\,2\mathrm{s}^2\,2\mathrm{p}^3$. Mg:
$1\mathrm{s}^2\,2\mathrm{s}^2\,2\mathrm{p}^6\,3\mathrm{s}^2$.
\end{solution}

\begin{solution}{exo:g10:electron-shells:2}
Carbono 4, nitrógeno 5, magnesio 2.
\end{solution}

\begin{solution}{exo:g10:electron-shells:3}
\omterm{def:g10:electron-shells:configuration}{Capa} externa $n = 3$: periodo 3. $2 + 3 = 5$ \omterm{def:g10:electron-shells:valence}{electrones de valencia}:
\omterm{def:g9:periodic-table-first-look:period}{grupo} 15 (fósforo).
\end{solution}

\begin{solution}{exo:g10:electron-shells:4}
\omterm{def:g10:electron-shells:family}{Metales alcalinos} (litio, sodio, potasio); \omterm{def:g10:electron-shells:family}{halógenos} (flúor, cloro,
bromo, yodo); \omterm{def:g10:electron-shells:family}{gases nobles} (helio, neón, argón).
\end{solution}

\begin{solution}{exo:g10:electron-shells:5}
Los \omterm{def:g7:atoms-and-molecules:atom}{átomos} tienden a ganar o perder \omterm{def:g9:inside-the-atom:nucleus}{electrones} para tener ocho
\omterm{def:g9:inside-the-atom:nucleus}{electrones} en su \omterm{def:g10:electron-shells:configuration}{capa} externa, como el \omterm{def:g10:electron-shells:family}{gas noble} más cercano.
\end{solution}

\begin{solution}{exo:g10:electron-shells:6}
K, $2.8.8.1$, pierde uno: \ce{K+},
$1\mathrm{s}^2\,2\mathrm{s}^2\,2\mathrm{p}^6\,3\mathrm{s}^2\,3\mathrm{p}^6$
(argón). F, $2.7$, gana uno: \ce{F-},
$1\mathrm{s}^2\,2\mathrm{s}^2\,2\mathrm{p}^6$ (neón).
\end{solution}

\begin{solution}{exo:g10:electron-shells:7}
$1\mathrm{s}^2\,2\mathrm{s}^2\,2\mathrm{p}^6\,3\mathrm{s}^2$: $Z = 12$,
magnesio.
\end{solution}

\begin{solution}{exo:g10:electron-shells:8}
\ce{Mg^{2+}} y \ce{Cl-}: \ce{MgCl2}. \ce{Al^{3+}} y \ce{O^{2-}}:
\ce{Al2O3}.
\end{solution}

\begin{solution}{exo:g10:electron-shells:9}
El azufre (6 \omterm{def:g10:electron-shells:valence}{electrones de valencia}, $2.8.6$), del mismo grupo, el 16.
\end{solution}

\begin{solution}{exo:g10:electron-shells:10}
Su \omterm{def:g10:electron-shells:configuration}{capa} externa ya está completa (ocho \omterm{def:g9:inside-the-atom:nucleus}{electrones}): no gana nada
perdiendo o ganando \omterm{def:g9:inside-the-atom:nucleus}{electrones}.
\end{solution}

\begin{solution}{exo:g10:electron-shells:11}
El argón tiene 18 \omterm{def:g9:inside-the-atom:nucleus}{electrones}; \ce{X^{2+}} ha perdido 2, así que X tiene
20: el calcio.
\end{solution}

\begin{solution}{exo:g10:electron-shells:12}
El helio, $1\mathrm{s}^2$, tiene una \omterm{def:g10:electron-shells:configuration}{capa} externa completa (la \omterm{def:g10:electron-shells:configuration}{capa} 1
solo admite dos \omterm{def:g9:inside-the-atom:nucleus}{electrones}), como los demás \omterm{def:g10:electron-shells:family}{gases nobles}; no se comporta
como los \omterm{def:g9:periodic-table-first-look:metal}{metales} del \omterm{def:g9:periodic-table-first-look:period}{grupo} 2.
\end{solution}

\begin{solution}{exo:g10:electron-shells:13}
Su \omterm{def:g10:electron-shells:valence}{electrón de valencia} está en la \omterm{def:g10:electron-shells:configuration}{capa} 4, más lejos del \omterm{def:g9:inside-the-atom:nucleus}{núcleo} que el
del sodio, en la \omterm{def:g10:electron-shells:configuration}{capa} 3: está menos fuertemente retenido y se pierde con
más facilidad.
\end{solution}

\begin{solution}{exo:g10:electron-shells:14}
Los cuatro tienen 18 \omterm{def:g9:inside-the-atom:nucleus}{electrones} y la configuración del argón:
\[ 1\mathrm{s}^2\,2\mathrm{s}^2\,2\mathrm{p}^6\,3\mathrm{s}^2\,3\mathrm{p}^6. \]
\end{solution}

\begin{solution}{exo:g10:electron-shells:15}
Gana 3 \omterm{def:g9:inside-the-atom:nucleus}{electrones} para completar el octeto: tiene 5 \omterm{def:g10:electron-shells:valence}{electrones de
valencia}, \omterm{def:g9:periodic-table-first-look:period}{grupo} 15. En el periodo 3, es el fósforo:
\[ 1\mathrm{s}^2\,2\mathrm{s}^2\,2\mathrm{p}^6\,3\mathrm{s}^2\,3\mathrm{p}^3. \]
\end{solution}

\begin{solution}{pb:g10:electron-shells:1}
\textbf{1.} Aluminio: 13 \omterm{def:g9:inside-the-atom:proton}{protones}, 13 \omterm{def:g9:inside-the-atom:nucleus}{electrones}. Oxígeno: 8 \omterm{def:g9:inside-the-atom:proton}{protones},
8 \omterm{def:g9:inside-the-atom:nucleus}{electrones}.

\textbf{2.} Al: $1\mathrm{s}^2\,2\mathrm{s}^2\,2\mathrm{p}^6\,3\mathrm{s}^2\,3\mathrm{p}^1$.
O: $1\mathrm{s}^2\,2\mathrm{s}^2\,2\mathrm{p}^4$.

\textbf{3.} Al: 3 \omterm{def:g10:electron-shells:valence}{electrones de valencia}, periodo 3, \omterm{def:g9:periodic-table-first-look:period}{grupo} 13. O: 6,
periodo 2, \omterm{def:g9:periodic-table-first-look:period}{grupo} 16.

\textbf{4.} El aluminio es un \omterm{def:g9:periodic-table-first-look:metal}{metal}; el oxígeno, un \omterm{def:g9:periodic-table-first-look:metal}{no metal}.

\textbf{5.} \ce{Al^{3+}}: $1\mathrm{s}^2\,2\mathrm{s}^2\,2\mathrm{p}^6$.

\textbf{6.} \ce{O^{2-}}: $1\mathrm{s}^2\,2\mathrm{s}^2\,2\mathrm{p}^6$.

\textbf{7.} El neón.

\textbf{8.} El aluminio pierde 3, y el oxígeno gana 2.

\textbf{9.} $2 \times (+3) + 3 \times (-2) = 0$: \ce{Al2O3}.

\textbf{10.} Perdidos: $2 \times 3 = 6$. Ganados: $3 \times 2 = 6$.

\textbf{11.} Los \omterm{def:g9:inside-the-atom:nucleus}{electrones} no se crean ni se destruyen: los que pierde
el aluminio son exactamente los que gana el oxígeno, y el compuesto es
neutro.

\textbf{12.} Tiene la misma carga, $+3$: las cargas del cristal siguen
compensándose.

\textbf{13.} $2 \times 27 + 3 \times 16 = 102$ \omterm{def:g9:inside-the-atom:proton}{nucleones}.

\textbf{14.} $102 \times 1.67 \times 10^{-27} = \qty{1.70e-25}{kg}$.

\textbf{15.} $10^{-3} / 1.70 \times 10^{-25} \approx 5.87 \times 10^{21}$
unidades fórmula.

\textbf{16.} $2 \times 5.87 \times 10^{21} = 1.17 \times 10^{22}$
\omterm{def:g9:ions:ion}{iones} \ce{Al^{3+}}; $3 \times 5.87 \times 10^{21} = 1.76 \times 10^{22}$
\omterm{def:g9:ions:ion}{iones} \ce{O^{2-}}.

\textbf{17.} Un \omterm{def:g9:inside-the-atom:nucleus}{electrón} es unas 1836 veces más ligero que un \omterm{def:g9:inside-the-atom:proton}{nucleón};
los 50 \omterm{def:g9:inside-the-atom:nucleus}{electrones} de una unidad fórmula pesan menos de tres centésimas
de un \omterm{def:g9:inside-the-atom:proton}{nucleón}.

\textbf{18.} $6 \times 5.87 \times 10^{21} \approx 3.5 \times 10^{22}$
\omterm{def:g9:inside-the-atom:nucleus}{electrones} transferidos para \qty{1}{g} de corindón.
\end{solution}
